\documentclass[UTF8,a4paper,11pt]{ctexart}
\usepackage[slantfont,boldfont]{xeCJK}
\usepackage{fontspec}

\usepackage{mathtools}
\usepackage{amsmath}
\usepackage{amsfonts}
\usepackage{amssymb}
\usepackage{amsthm}
\usepackage{indentfirst}
\usepackage{graphicx}
\usepackage{geometry}
\usepackage{fancyhdr}
\usepackage{titlesec}
\usepackage{setspace}
\usepackage{enumitem}
\usepackage{bm}
\usepackage{hyperref}

\setlength{\parindent}{2em}
\linespread{1.2}
\geometry{left=2.5cm,right=2.5cm,top=2.5cm,bottom=2.5cm}

\everymath{\displaystyle}

\newcommand*{\QED}{\hfill\ensuremath{\square}}
\newcommand{\inner}[2]{\left(#1,#2\right)}

\setlength{\headheight}{13.6pt}
\pagestyle{fancy}
\fancyhf{}
\fancyfoot[C]{\thepage}
\fancyhead[L]{数值分析\quad 第五次理论作业\quad 学号：241098117，姓名：张城宗}
\fancyhead[R]{习题三}

\begin{document}

\begin{center}
  {\LARGE\bfseries 数值分析\quad 第五次理论作业}\\[6pt]
  {\large 习题三第 7、8、9 题}\\[4pt]
  {\normalsize 学号：241098117\quad 姓名：张城宗}
\end{center}

\bigskip

%=====================================================
\noindent\textbf{习题三\quad 第 7 题}\quad
对下列给定的权函数 $W(x)$，求区间 $[-1,1]$ 上的正交多项式系的前三个多项式
$p_0(x),p_1(x),p_2(x)$：
\[
  (1)\ W(x)=\frac{1}{\sqrt{1+x^2}},\qquad
  (2)\ W(x)=1+x^2.
\]
\medskip
\noindent\textbf{解}\quad
由讲义第三章性质 3.3.2，区间 $[-1,1]$ 上关于权函数 $W(x)$ 的首一正交多项式满足
\[
  p_{k+1}(x)=(x-\alpha_k)p_k(x)-\beta_k p_{k-1}(x),\qquad
  p_0(x)=1,\quad p_{-1}(x)=0,
\]
其中
\[
  \alpha_k=\frac{\inner{x p_k}{p_k}}{\inner{p_k}{p_k}},\qquad
  \beta_1=\frac{\inner{p_1}{p_1}}{\inner{p_0}{p_0}},
\]
这里的内积为
\[
  \inner{f}{g}=\int_{-1}^1 W(x)f(x)g(x)\,\mathrm{d}x.
\]
由于两种情形下的权函数都满足 $W(-x)=W(x)$，所以它们都是偶函数，从而
\[
  \alpha_0=\frac{\int_{-1}^1 xW(x)\,\mathrm{d}x}{\int_{-1}^1 W(x)\,\mathrm{d}x}=0,
\]
故
\[
  p_0(x)=1,\qquad p_1(x)=x.
\]
同理，
\[
  \alpha_1=\frac{\int_{-1}^1 x^3W(x)\,\mathrm{d}x}{\int_{-1}^1 x^2W(x)\,\mathrm{d}x}=0,
\]
所以
\[
  p_2(x)=x^2-\beta_1,\qquad
  \beta_1=\frac{\int_{-1}^1 x^2W(x)\,\mathrm{d}x}{\int_{-1}^1 W(x)\,\mathrm{d}x}.
\]
下面分别计算 $\beta_1$。

\medskip
\noindent\textbf{(1) 当 $W(x)=\dfrac{1}{\sqrt{1+x^2}}$ 时}\quad
\[
  \inner{p_0}{p_0}=\int_{-1}^1 \frac{1}{\sqrt{1+x^2}}\,\mathrm{d}x
  =2\int_0^1 \frac{1}{\sqrt{1+x^2}}\,\mathrm{d}x
  =2\,\operatorname{arsinh}1
  =2\ln(1+\sqrt{2}).
\]
再计算
\[
  \inner{p_1}{p_1}=\int_{-1}^1 \frac{x^2}{\sqrt{1+x^2}}\,\mathrm{d}x
  =2\int_0^1 \frac{x^2}{\sqrt{1+x^2}}\,\mathrm{d}x.
\]
令 $x=\sinh t$，则 $\mathrm{d}x=\cosh t\,\mathrm{d}t$，$\sqrt{1+x^2}=\cosh t$，于是
\begin{align*}
  \inner{p_1}{p_1}
  &=2\int_0^{\operatorname{arsinh}1}\sinh^2 t\,\mathrm{d}t \\
  &=2\left[\frac{\sinh 2t}{4}-\frac{t}{2}\right]_0^{\operatorname{arsinh}1} \\
  &=2\left(\frac{2\sqrt{2}}{4}-\frac{1}{2}\ln(1+\sqrt{2})\right) \\
  &=\sqrt{2}-\ln(1+\sqrt{2}).
\end{align*}
故
\[
  \beta_1=\frac{\sqrt{2}-\ln(1+\sqrt{2})}{2\ln(1+\sqrt{2})}.
\]
因此前三个首一正交多项式为
\[
  \boxed{
  p_0(x)=1,\qquad
  p_1(x)=x,\qquad
  p_2(x)=x^2-\frac{\sqrt{2}-\ln(1+\sqrt{2})}{2\ln(1+\sqrt{2})}.}
\]

\medskip
\noindent\textbf{(2) 当 $W(x)=1+x^2$ 时}\quad
\[
  \inner{p_0}{p_0}=\int_{-1}^1 (1+x^2)\,\mathrm{d}x
  =2+\frac{2}{3}
  =\frac{8}{3},
\]
\[
  \inner{p_1}{p_1}=\int_{-1}^1 x^2(1+x^2)\,\mathrm{d}x
  =\int_{-1}^1 (x^2+x^4)\,\mathrm{d}x
  =\frac{2}{3}+\frac{2}{5}
  =\frac{16}{15}.
\]
于是
\[
  \beta_1=\frac{16/15}{8/3}=\frac{2}{5}.
\]
故前三个首一正交多项式为
\[
  \boxed{
  p_0(x)=1,\qquad
  p_1(x)=x,\qquad
  p_2(x)=x^2-\frac{2}{5}.}
\]

\bigskip\bigskip

%=====================================================
\noindent\textbf{习题三\quad 第 8 题}\quad
设 $q_n(x)$ 是区间 $[a,b]$ 上关于权函数 $W(x)$ 的首一 $n$ 次正交多项式。令
\[
  A_n=
  \begin{pmatrix}
    \alpha_0 & \sqrt{\beta_1} &        &        &   \\
    \sqrt{\beta_1} & \alpha_1 & \sqrt{\beta_2} &        &   \\
        & \sqrt{\beta_2} & \ddots & \ddots &   \\
        &        & \ddots & \alpha_{n-2} & \sqrt{\beta_{n-1}} \\
        &        &        & \sqrt{\beta_{n-1}} & \alpha_{n-1}
  \end{pmatrix},
  \qquad n\ge 1,
\]
其中 $\alpha_k,\beta_k$ 为递推关系式 $(3.3.3)$ 中的系数。试证，矩阵 $A_n$ 的特征值是正交多项式
$q_n(x)$ 的根。
\medskip
\noindent\textbf{证明}\quad
由讲义第三章性质 3.3.2，首一正交多项式满足三项递推关系
\[
  q_{k+1}(x)=(x-\alpha_k)q_k(x)-\beta_k q_{k-1}(x),
  \qquad q_0(x)=1,\quad q_1(x)=x-\alpha_0.
\]
记 $\Delta_n(\lambda)=\det(\lambda I-A_n)$ 为矩阵 $A_n$ 的特征多项式。再设
$A_k$ 表示 $A_n$ 的左上角 $k\times k$ 主子矩阵，并约定
\[
  \Delta_0(\lambda)=1,\qquad \Delta_1(\lambda)=\lambda-\alpha_0.
\]

对 $k\ge 2$，按最后一行或最后一列展开 $\Delta_k(\lambda)=\det(\lambda I-A_k)$，得
\[
  \Delta_k(\lambda)
  =(\lambda-\alpha_{k-1})\Delta_{k-1}(\lambda)-\beta_{k-1}\Delta_{k-2}(\lambda).
\]
这是因为相邻副对角元都等于 $\sqrt{\beta_{k-1}}$，其乘积恰为 $\beta_{k-1}$。

将正交多项式的递推式把指标改写为 $k\mapsto k-1$，可得
\[
  q_k(x)=(x-\alpha_{k-1})q_{k-1}(x)-\beta_{k-1}q_{k-2}(x),
  \qquad k\ge 2.
\]
于是 $\Delta_k(\lambda)$ 与 $q_k(\lambda)$ 满足完全相同的递推关系与初值条件：
\[
  \Delta_0(\lambda)=1=q_0(\lambda),\qquad
  \Delta_1(\lambda)=\lambda-\alpha_0=q_1(\lambda).
\]
因此由递推关系的唯一性可知，
\[
  \Delta_k(\lambda)=q_k(\lambda),\qquad k\ge 0.
\]
特别地，
\[
  \det(\lambda I-A_n)=\Delta_n(\lambda)=q_n(\lambda).
\]
故矩阵 $A_n$ 的特征值满足
\[
  \det(\lambda I-A_n)=0,
\]
也就是正交多项式 $q_n(\lambda)=0$ 的根。证毕。\QED

\bigskip\bigskip

%=====================================================
\noindent\textbf{习题三\quad 第 9 题}\quad
设
\[
  X_n(x,y)=T_{n+1}(x)T_n(y)-T_{n+1}(y)T_n(x),
\]
其中 $T_n$ 表示 Chebyshev 多项式。证明
\[
  X_n(x,y)=2(x-y)T_n(x)T_n(y)+X_{n-1}(x,y),\qquad n\ge 1,
\]
并导出
\[
  \frac{1}{2}X_n(x,y)=(x-y)\left[\frac{1}{2}+\sum_{k=1}^n T_k(x)T_k(y)\right].
\]
\medskip
\noindent\textbf{证明}\quad
由讲义第三章中 Chebyshev 多项式的递推关系
\[
  T_{n+1}(t)=2tT_n(t)-T_{n-1}(t),\qquad n\ge 1,
\]
分别在 $t=x$ 与 $t=y$ 处使用，代入 $X_n(x,y)$ 得
\begin{align*}
  X_n(x,y)
  &=[2xT_n(x)-T_{n-1}(x)]T_n(y)-[2yT_n(y)-T_{n-1}(y)]T_n(x) \\
  &=2(x-y)T_n(x)T_n(y)+T_n(x)T_{n-1}(y)-T_n(y)T_{n-1}(x) \\
  &=2(x-y)T_n(x)T_n(y)+X_{n-1}(x,y).
\end{align*}
故第一式成立。

\medskip
由上式移项得
\[
  X_n(x,y)-X_{n-1}(x,y)=2(x-y)T_n(x)T_n(y),\qquad n\ge 1.
\]
从 $k=1$ 到 $n$ 求和，得到
\[
  X_n(x,y)-X_0(x,y)=2(x-y)\sum_{k=1}^n T_k(x)T_k(y).
\]
又因为
\[
  X_0(x,y)=T_1(x)T_0(y)-T_1(y)T_0(x)=x-y,
\]
其中 $T_0(x)=1,\ T_1(x)=x$，所以
\[
  X_n(x,y)=(x-y)+2(x-y)\sum_{k=1}^n T_k(x)T_k(y).
\]
提取公因子 $(x-y)$，即得
\[
  X_n(x,y)=2(x-y)\left[\frac{1}{2}+\sum_{k=1}^n T_k(x)T_k(y)\right].
\]
两边同除以 $2$，便有
\[
  \boxed{
  \frac{1}{2}X_n(x,y)=(x-y)\left[\frac{1}{2}+\sum_{k=1}^n T_k(x)T_k(y)\right].}
  \qquad \QED
\]

\end{document}
